\section{Methodology}

A vertical pipeline, like the one proposed in this paper, is endogenous. The `first layer' is the raw text processed through a pre-trained natural language inference model (DeBERTa-v2-xlarge-mnli) in which we introduce millions of tiny, unavoidable pre-justified premises. Initially, this architecture was designed to help mathematically formalize the previously mentioned relativism in an attempt to demonstrate the myth of objective neutrality. By concentrating different perspectives with adversarial framing, however, it became clear that the system ultimately functions as an attempt at measuring the ideal of an Overton Window (what is perceived as the `middle ground' or the range of politically acceptable ideas) through the mapping of structural invariants and relativistic variations. By treating bias as a physical shape with physical properties, we can perform various transformations that help us make certain inferences about the structure of semantics in the context of language models.

\subsection{Adversarial Framing and the V-Observers}

The proposed system takes one frozen NLI bot and gives it eight different adversarial or contrastive hypotheses, prompting it to analyze any given number of articles. In this paper, we refer to these different conditioned perspectives as `V-Observers'. In our implementation we used both real and synthetic articles relating to the Palestinian genocide. The mechanism of adversarial framing allows our approach to dialectically process each article under paired framing hypotheses rather than the raw average embeddings. 

The system queries a series of hypothesis pairs designed to stress-test the narrative structure:

\begin{itemize}
    \item \textbf{Pair 1: Israeli actions framing} \\
    A: ``This article portrays Israeli military actions as legitimate self-defense against terrorism.'' \\
    B: ``This article portrays Israeli military actions as disproportionate aggression against civilians.''
    \item \textbf{Pair 2: Palestinian framing} \\
    A: ``This article portrays Palestinians as innocent victims of occupation and violence.'' \\
    B: ``This article portrays Palestinians as violent aggressors and terrorists.''
    \item \textbf{Pair 3: Proportionality framing} \\
    A: ``This article describes Israeli responses as proportionate and measured.'' \\
    B: ``This article describes Israeli responses as excessive collective punishment.''
    \item \textbf{Pair 4: Resistance vs terrorism framing} \\
    A: ``This article frames Palestinian armed groups as legitimate resistance movements.'' \\
    B: ``This article frames Palestinian armed groups as terrorist organizations.''
    \item \textbf{Pair 5: Occupation framing} \\
    A: ``This article frames the situation as Israeli occupation of Palestinian territory.'' \\
    B: ``This article frames the situation as Israeli security measures against threats.''
    \item \textbf{Pair 6: Civilian casualties framing} \\
    A: ``This article suggests Israeli forces deliberately target civilians.'' \\
    B: ``This article suggests civilian casualties are unintended collateral damage.''
    \item \textbf{Pair 7: Historical context framing} \\
    A: ``This article frames Israel as a colonial settler state.'' \\
    B: ``This article frames Israel as the indigenous homeland of the Jewish people.''
    \item \textbf{Pair 8: International law framing} \\
    A: ``This article emphasizes Israeli violations of international law.'' \\
    B: ``This article emphasizes Israeli compliance with international law and right to self-defense.''
\end{itemize}

By forcing the text to pass through these specific dialectical gates we can measure the friction generated when a narrative is stretched between two opposing coordinated frameworks of reality.

\subsection{Track 1: Logits and the Substrate}

The first step of the pipeline is ingestion. We take a target article and run it through the eight adversarial NLI gates. The model outputs a three-dimensional score: entailment, neutral, and contradiction. In implementation, this becomes a raw 24-dimensional verdict channel. These logits function strictly as the baseline diagnostic of the model's forward-pass verdict behavior. Track 1 is the constitutional floor of the system: the raw verdict substrate against which later geometric transformations can be compared.

\subsection{Track 1.5: Spectral Polarity Extraction}

Logits do not provide specifics as to why the model made a decision. Track 1.5 measures the article's weighted resistance to reinterpretation by extracting the semantic shift induced by contrastive observer frames. The system extracts the observer representations generated by the paired hypotheses and computes the A-minus-B polarity deltas across the eight contrastive pairs. These deltas form an eight-way aggregate field of local semantic shear. 

The ``Polarity Delta'' measures how the article's semantic position shifts when it is contrasted against opposing observer frames. Where the field is steep, the article is resistant to reinterpretation; where it is flatter, the article has more leeway of interpretation. By performing Singular Value Decomposition on the centered contrastive delta matrices, the system endogenously discovers the primary axes of semantic stress without human intuition dictating which side of the conflict must dominate. Track 1.5 is the system's compass: it discovers the dominant direction of semantic shear but does not decide what that shear morally means.

\subsection{Track 2: Geometric Expansion and the JL Lemma}

Track 2 supplies the coordinate container necessary to position Track 1.5 shear in relation to a latent topology. To formalize this, we invoke the Johnson-Lindenstrauss (JL) lemma, which mathematically justifies why the semantic structure is robust across random projections but perspective-dependent. Unlike logits, embeddings have no automatic shared currency. The solution is to enforce a shared spatial geometry, making the observer views born-aligned in the same coordinate system using Random Kitchen Sinks \cite{Rahimi2008} to approximate a kernel feature space.

The kernel-approximated manifold allows observer shear to warp the effective distance between articles. Conceptually, we write the Observer Metric Tensor as:
\begin{equation}
g_{\mu\nu}^{total}(x) = \frac{1}{\rho(x)}\delta_{\mu\nu} + \nabla_\mu \Phi \nabla_\nu \Phi
\end{equation}
where $\delta_{\mu\nu}$ represents the base content-to-content geometry, $\nabla_\mu \Phi \nabla_\nu \Phi$ represents semantic shear from Track 1.5, and $\rho(x)$ is the conformal density from Track 3.

\subsection{Track 3: The Dirichlet Sweep and Homotopy Type Theory}

Track 3 provides a lens. Inside the model, there is no guaranteed continuous plane where the distance between every two points is meaningful. Formalized through Homotopy Type Theory (HoTT), Track 3 asks whether paths (homotopies) between semantic identities remain stable under continuous deformation of the observer frame. 

Mechanically, Track 3 uses a dynamically optimized weight simplex over RKHS-projected observer views. The system begins in a high-alpha consensus state with strong prior regularization, then cools toward lower alpha values where structured observer separation can emerge. The annealing objective balances separation and inertia:
\begin{equation}
L_t(w_t) = -\alpha_t^{-1}\,\mathrm{Var}\left(\sum_b w_{t,b}\phi_b(x)\right) + \lambda_{KL}\alpha_t\,\mathrm{KL}(w_t||w_{t-1})
\end{equation}
Track 3 produces a conformal density field $\rho_i$. Articles that maintain their observer-mixture position under annealing (stable homotopies) produce denser, more stable local geometry. Articles that tear apart produce lower-density regions.

\subsection{Track 4: Traversal Telemetry}

Track 4 simulates particles across the article graph based on density, stress, and metric distance to measure the work required to make semantic leaps (geodesics) between article nodes. It is a thermodynamic stochastic walk over a metric-respecting article graph. For each possible neighbor $j$, the transition weight is:
\begin{equation}
w_{ij} = \exp\left(-\frac{d_g(x_i,x_j)}{\tau}\right)
\end{equation}
where $d_g$ is the discrete metric edge cost induced by the fused geometry. When accumulated work reaches the cognitive horizon, the walker flips into retreat mode penalized toward the anchor coordinate. Track 4 exports mechanical telemetry: path coordinates, work integrals, and survival booleans, allowing us to study the structural traversability of the discourse manifold.

\subsection{Track 5: Synthesis and Procrustes Alignment}

Track 5 is the synthesis layer, functioning as the interferometer proper. It relies on kernel and metric fusion. In the Hadamard branch, Track 2 and Track 1.5 are converted into similarity kernels and multiplied element-wise:
\begin{equation}
K_{fused} = K_{T2} \odot K_{T1.5}
\end{equation}
Track 3 is then used as a conformal density normalizer to stretch or compress the fused manifold based on consensus and disagreement. This creates a unified representation where Procrustes alignment can be performed across contrasting ideological spaces to measure residual divergence.
